\begin{resumo}
Este trabalho busca propor uma análise comparativa entre dois métodos de processamento
para sistemas de \gls{smart_metering} (medição inteligente): \gls{edge_computing} (Computação em Borda), 
com inferência de \gls{ocr} realizada localmente no microcontrolador \gls{esp32cam} 
por meio de modelo de aprendizado de máquina treinado de forma supervisionada 
e quantizado em inteiros de 8 bits, 
e \gls{cloud_computing} (Computação em Nuvem), com inferência realizada em ambiente 
remoto através de modelo treinado de forma supervisionada e quantizado em ponto flutuante de 32 bits.
A motivação reside no elevado custo de substituição dos mais diversos medidores convencionais, 
amplamente utilizados em ambientes residenciais e privados, 
por modelos inteligentes nativos, o que favorece o \textit{\gls{retrofitting}} 
como alternativa mais economicamente viável. 
A transmissão dos dados entre os nós de borda e o gateway central é realizada via rede \gls{lorawan},
utilizando o Spreading Factor SF8, configuração que assegura conformidade com o limite
regulatório de \textit{dwell time} de 400~ms estabelecido pelo Ato n\textordmasculine~14.448
da ANATEL, ao mesmo tempo em que maximiza o alcance dentro desse limite. A viabilidade
do envio de imagens fragmentadas via \gls{lorawan} foi demonstrada analiticamente, com
uma imagem de 160$\times$120 pixels em escala de cinza, comprimida utilizando o algoritmo JPEG com
aproximadamente 3,10~KB, podendo ser transmitida em 38 pacotes em cerca de 13,3~s
de rádio ativo. Os dois fluxos de processamento
serão avaliados segundo três métricas quantitativas: acurácia de leitura \gls{ocr} (\%),
latência de ponta a ponta (ms) e consumo energético por leitura (mAh). 
O objetivo deste trabalho é apresentar resultados que orientem a escolha arquitetural
mais adequada para implantações de \textit{Smart Metering} em larga escala, considerando
o \textit{trade-off} entre eficiência energética, precisão de reconhecimento e complexidade
de manutenção do sistema.

\vspace{\onelineskip}

\noindent
\textbf{Palavras-chave}: Smart Metering. ESP32. LoRaWan. OCR.
\end{resumo}
